% Section fragment. Packages: amsmath, amssymb, amsthm, booktabs, hyperref.
\section{PA194GH: explicit-binder AX with cached-head Shift}
\label{sec:pa194gh-new}

PA194GH combines the supplied G emitter with the supplied H controller.
Its saved generalization uses the explicit AX operand, and its Shift
tests only the cached theorem. Neither ingredient is new in isolation.
Their combination changes the command compiler and is absent from all
94 recognized 194-state transition graphs at the pinned snapshot.
The construction has exactly 194 working states and a separate halt label.
It is a new 194-state fallback, not a state-count improvement.

The unconditional results are the literal table, its finite construction
checks, the symbolic unbounded-register contracts, and the command-word
identities below. At the level of ordinary finite computations these give
the same refutation-search behavior as the numerical calculus
$\mathcal D_{4,Q}$. The absolute iPRA equivalence with PA consistency
still depends on the missing internal calibration. This section does not
treat Python certificates as proof-assistant theorems.

\subsection{Exact numerical interface and relation to G}

Let $p(a,b)=(a+b)(a+b+1)/2+b$, $[a;w]=p(a,w)$, and use total Cantor
projections. Thus $[0;0]=0$ and $\operatorname{hd}(0)=\operatorname{tl}(0)=0$.
The constructors are
\[
 I(a,b)=4p(a,b),\quad M(a,b)=4p(a,b)+1,\quad
 U(i,a)=4p(a,i)+2,\quad E(a,b)=4p(a,b)+3.
\]
Code zero is truth, and one is the refutation target. Write
$a\to b=I(a,b)$, $N(a)=I(a,1)$,
$a\wedge b=N(I(a,N(b)))$, and $\exists_i a=N(U(i,N(a)))$.
The exact heads inherited from G are
\begin{align*}
 h_1(D,H)&=(\exists_D U(D,H))\to H,\\
 h_2(D,A,H)&=U(D,U(A,H))\to U(A,U(D,H)),\\
 h_3(D,A,H)&=M(H,D)\to(E(H,A)\to M(A,D)),\\
 h_4(D,A,H)&=E(H,D)\to(M(A,H)\to M(A,D)),\\
 h_5(D,H)&=\exists_D E(H,D),\\
 h_6^Q(D,A,H)&=U(D,A\to H)\to(U(D,A)\to U(D,H)),\\
 \rho(D,A,H,u)&=((H\to u)\to D)\to((D\to H)\to(A\to H)),
                  \qquad u\le D,\\
 \Gamma(A)&=A\to U(A,A).
\end{align*}
For completeness, the inherited gate is defined by
\begin{align*}
 P(a,b,c)&=c\to M(c,b),\qquad
 V(a,b,c)=M(a,b)\to U(c,M(a,b)),\\
 S_0(a,b,c)&=U\bigl(b,(M(a,b)\to P(a,b,c))\wedge(M(c,b)\to c)\bigr),\\
 Q_0(a,b,c)&=U\bigl(c,S_0(a,b,c)\to N(P(a,b,c))\bigr),\\
 g(a,b,c)&=\begin{cases}
 N(V(a,b,c)\to Q_0(a,b,c)),&c\ne a\text{ and }c\ne b,\\
 0,&\text{otherwise}.
 \end{cases}
\end{align*}
The calculus has these axiom instances, truth, modus ponens, and unrestricted
generalization of theorems. The Q head is the supplied exchange of the two
curried premises of universal distribution. No further logical rule is
introduced by GH.

A store is $(t,w,d,h)$: visible theorem, stack, numerical offset, cached
theorem. For $D=d+x$ and $H=h+y$, AX is \emph{exactly G's command}:
\[
 \mathsf{AX}_{c,x,y}(t,w,d,h)
       =(g(H,D,c),[q_1,q_2,\ldots,q_9,t;w],0,0),
\]
where the top-first generated entries are
\[
 (q_1,\ldots,q_9)=
 (\Gamma(c),h_6^Q(D,c,H),h_5(D,H),h_4(D,c,H),
 h_3(D,c,H),h_2(D,c,H),h_1(D,H),\rho(D,c,H,d),U(c,h)).
\]
In particular $q_9$ uses the raw incoming $h$, not $H$, and the explicit
binder $c$, not the offset $d$. The entire emitter word and its complete
output stack are unchanged from G and checked in a free term algebra.
All of $g,q_1,\ldots,q_8$ are zero modulo four; $q_9$ is two modulo four
and at least two. These observations do not assume consistency.

Writing $v=d+x$, MP is G's successful no-push command. Shift is H's:
\begin{align*}
 \mathsf{MP}_x(t,w,d,h)&=
 \begin{cases}(v,w,v,0),&I(t,v)=h,\\(0,w,v,0),&I(t,v)\ne h;\end{cases}\\
 \mathsf S(t,w,d,h)&=
 \begin{cases}\mathsf{halt},&h=1,\\
 (h,\operatorname{tl}(w),d,\operatorname{hd}(w)),&h\ne1.
 \end{cases}
\end{align*}
The change is observable: Shift from $(1,w,0,0)$ proceeds in GH but halts
in G. Conversely, on $(1,w,0,1)$ it halts in GH but G's sum guard is false.
At $d=h=0,c=1$, GH's saved generalization is $U(1,0)=10$, whereas H's is
$U(0,0)=2$. Thus GH differs from both parents at the command interface.

\subsection{Uniform words, including zero aliases}

Commands execute from left to right. For every store,
\begin{align*}
 \mathsf{AX}_{c,x,y};\mathsf S^9(t,w,d,h)
     &=(q_8,[t;w],0,U(c,h)),\\
 \mathsf{AX}_{c,x,y};\mathsf S^{10}(t,w,d,h)
     &=(U(c,h),w,0,t).
\end{align*}
These Shifts test only the cleared cache or generated heads, never the
arbitrary saved top. Put $A_0=\mathsf{AX}_{0,0,0}$ and define
\[
 \mathsf{Push}=A_0;\mathsf S^9;\mathsf{MP}_0,
 \quad \mathsf{Pop}=A_0;\mathsf S^{11},
 \quad \mathsf{Discard}=\mathsf{MP}_0.
\]
The exact contracts are
\begin{align*}
 \mathsf{Push}(t,w,d,h)&=(0,[t;w],0,0),\\
 \mathsf{Pop}(t,[a;w],d,h)&=
 \begin{cases}\mathsf{halt},&t=1,\\(t,w,0,a),&t\ne1;\end{cases}\\
 \mathsf{Discard}(t,w,0,h)&=(0,w,0,0).
\end{align*}
Push's final comparison fails because implication and universal codes have
different residues. Pop's eleventh Shift tests exactly the saved $t$; it
is its only possible halt. Discard has candidate conclusion zero, so its
two branches coincide for every old cache. The equations also hold for
the zero stack using the total projections: in particular Push does not
create a distinguishable zero entry when both $t$ and $w$ are zero.

Let $\mathsf{Offset}_i=\mathsf{MP}_i$. From $(0,w,0,0)$ it returns
$(0,w,i,0)$, including at $i=0$. For $1\le j\le7$ the selector
\[
 \mathsf{Discard};\mathsf{AX}_{c,D,H};\mathsf S^{j+1};
             \mathsf{Pop}^{9-j}
\]
returns $(q_j,w,0,0)$ from every $(t,w,0,0)$. The selected head is never
one, so every Pop is safe. The final consumed entry is the saved zero.
The gate selector is $\mathsf{Discard};\mathsf{AX}_{c,D,H};\mathsf{Pop}^{10}$.
For $u\le D$, the selector
\[
 \mathsf{Discard};\mathsf{Offset}_u;
       \mathsf{AX}_{c,D-u,H};\mathsf S^9;\mathsf{Pop}
\]
returns $(\rho(D,c,H,u),w,0,0)$. Truth uses Discard.

G's shorter generalizer survives the change to Shift:
\[
 \mathsf{Gen}_i=\mathsf{Push};\mathsf S;
                    \mathsf{AX}_{i,0,0};\mathsf S^{10}.
\]
From $(a,w,0,0)$, the first two operations give $(0,w,0,a)$; AX then
consumes the cache and saves $U(i,a)$, and ten Shifts return
$(U(i,a),w,0,0)$. This execution never halts, including when $a=1$.
No offset-setting command is needed because the binder comes from the
explicit AX operand.

\subsection{All-store refutation compiler and exhaustive scheduler}

The framed MP suffix has the exact alternative
\[
 (\mathsf S;\mathsf S;\mathsf{MP}_b;\mathsf{Pop})
         (0,[a,I(a,b),0;w],0,0)
  =\begin{cases}
       \mathsf{halt},&a=1\text{ or }b=1,\\
       (b,w,0,0),&a\ne1\text{ and }b\ne1.
    \end{cases}
\]
The first Shift caches $a$. The second either halts or produces
$(a,[0;w],0,I(a,b))$. Successful MP gives $(b,[0;w],b,0)$, after which
Pop either detects $b=1$ or consumes the reserved zero and clears the
offset. This includes $a=b=w=0$, since both the stack and implication
zero aliases are identities of the numerical constructors.

Compile a finite $\mathcal D_{4,Q}$ proof tree $e$ into $C(e)$ as follows.
Use the selectors for axiom nodes and append $\mathsf{Gen}_i$ for a
generalization node. For an MP node with minor proof $e_0$ of $a$ and major
proof $e_1$ of $I(a,b)$, use
\[
 \mathsf{Discard};\mathsf{Push};C(e_1);\mathsf{Push};
 C(e_0);\mathsf{Push};\mathsf S;\mathsf S;\mathsf{MP}_b;\mathsf{Pop}.
\]
Induction on the proof tree shows that from any $(t,w,0,0)$ this word either
halts or returns $(\operatorname{concl}(e),w,0,0)$. The initial Discard and
Push reserve the zero required by the MP suffix. In the generalization case
the induction hypothesis either supplies a halt or the exact input to the
total shorter generalizer. If $e$ proves one, then
\[
                         A_0;C(e);\mathsf{Pop}
\]
halts from every store, including stores with arbitrary offset and cache.
The first AX clears those fields, and final Pop detects one even when it
was not produced by an MP node. A final Shift alone would be incorrect.

The scheduler is unchanged. Its command and continuation codes are
\begin{align*}
 B_{\mathsf S}(z)&=p(0,z),\qquad
 B_{\mathsf{MP}_x}(z)=p(1,p(x,z)),\\
 B_{\mathsf{AX}_{c,x,y}}(z)&=p(c+2,p(y,p(x,z))),\\
 E(\epsilon)&=0,\qquad E(c::w)=1+B_c(E(w)).
\end{align*}
For $z=E(A_0;C(e);\mathsf{Pop})$, prefixing this universal halting word
by $\mathsf{AX}_{k,0,0}$ gives fresh-candidate codes
\[
                       n_k=p(k+2,p(0,p(0,z))).
\]
They are strictly increasing and unbounded. The carried store after the
prefix is harmless by the all-store property. Positive continuation codes
are decoded after subtracting one and have strictly decreasing residuals.
Each prior candidate therefore finishes or halts, and a finite search
reaches a suitable $n_k$ unless it has halted earlier. The leading-one
distinction between a continuation code and a fresh candidate is essential.

For the converse, annotate the visible theorem, cached theorem, and each
live canonical stack entry with a $\mathcal D_{4,Q}$ derivation. The virtual
zero tail has the truth derivation. No annotation is required for the offset.
AX adds the listed axioms and the generalization $U(c,h)$ of an annotated
theorem. MP composes proofs on success and writes truth on failure. Shift
transfers annotations and halts only on a cached derivation of one.
This supplies the finite syntactic equivalence between scheduler halting
and $\mathcal D_{4,Q}$ refutability.

\subsection{Literal source and physical verification}

In the frozen G source DAG only the following named definition is replaced:
\[
 \texttt{node170}:=
       ((\texttt{node36},(\texttt{node167},\texttt{node145})),\texttt{node57}).
\]
The existing no-push MP definition
$\texttt{node163}=(\texttt{node6},\texttt{node54})$ remains.
The Shift branch now clears the old top, copies the cache into the top and
comparison scratch register, restores the cache, and enters the existing
destructive equality guard with scratch value one. Its guard is therefore
$h=1$. The full source contains 3,653 primitive sites and 4,754 outcomes;
the same G AX word has 190 macro tokens. The five native transition edits
relative to the G raw table are
\[
\begin{array}{cc|ccc|ccc}
q&b&\multicolumn{3}{c|}{\text{old write, move, target}}&
      \multicolumn{3}{c}{\text{new write, move, target}}\\
38&0&0&R&28&0&R&11\\
38&1&1&R&20&1&R&122\\
214&0&0&R&188&0&R&312\\
214&1&1&R&312&1&R&20\\
312&0&0&R&11&0&R&188
\end{array}
\]
They are the supplied H controller changes, applied to the G association.
The arithmetic register-engine rows are unchanged.

There are 134 fixed controls, 67 reader controls, and seven simultaneous
contextual folds, giving $134+67-7=194$. All 28 local fold cases close.
The construction checker verifies all 4,754 symbolic outcomes using
1,512,536 symbolic steps, with maximum 479 per outcome. An independent
literal-row checker rederives four register-list induction lemmas, uses two
independent source resolvers, and checks the same outcomes in 1,102,232
symbolic steps, maximum 385, with minimum positive duration constant three.
Its domain contains nine arbitrary natural registers, either decrement
branch, an arbitrary nonempty finite list of natural spare values, an
arbitrary nonnegative affine header gap, and complete blank exteriors.
These are symbolic contracts rather than finite-value sampling.

The complete finite-support blank tape reaches native dispatcher 84 at
source site 25, head $-2$, logical registers $(2,0,\ldots,0)$, and 26 spare
blocks of value one in exactly 57,922 transitions. Here the nonzero register
is the fresh-candidate counter; the command store is initially clear.
The 388-row dense table has start zero, working states $0,\ldots,193$,
and halt label 194. Its SHA-256 is
\begin{quote}\small
\nolinkurl{b9fadcb197c33895603c71b7119bc92b1772c01960804e8850e2a66dd39a995e}.
\end{quote}

Exact structural comparison against the entire recognized 94-graph frontier
finds no match under control renaming, choice of spanning root, reflection,
global symbol exchange, or ignored final halt effects. Symbol exchange is
used only as a stronger structural novelty test, not as an assertion of
zero-blank semantic equivalence.

\subsection{Arithmetic conclusion and remaining obligation}

The source-word compiler, annotation update, scheduler witnesses, and
finite physical simulations are explicit primitive-recursive operations.
Ordinary mathematical induction proves their finite-trace claims. These
claims support the usual ordinary-PA argument when composed with the
inherited numerical-calculus-to-PA refutation calibration. That calibration
is an explicit inherited dependency, not a consequence of this construction's
runtime certificates.

For an iPRA conclusion the required premises must themselves be established
in iPRA: (i) the exact source/tape trace reductions and the syntactic
soundness/completeness reductions just described, and (ii) the equivalence
of $\mathcal D_{4,Q}$ refutability with PA refutability. Once these premises
are supplied, constructive composition gives
\[
 \mathrm{iPRA}\vdash
 \neg\operatorname{Halt}_{\rm blank}(\mathrm{PA194GH})
       \leftrightarrow\operatorname{Con}(\mathrm{PA}).
\]
Neither the executable certificates nor a citation of the generic
\texttt{Calibrated} theorem supplies the missing premises. The current
repository reports the machine-specific iPRA calibration as unfinished.
Accordingly, GH is delivered as a fully explicit new 194-state design with
checked computational interfaces and an explicit conditional calibration,
not as a completed unconditional solution of that formalization problem.
